%=====================================================================
%  ANALYSIS OF ALGORITHMS --- FALL 2026
%  One-page weekly lecture schedule
%
%  SOURCE OF RECORD: parts/G_schedule.tex in the handout.  Change the
%  handout first, then this page -- the slide pipeline parses the
%  appendix, not this file.  Follows HEC 2023 (Curriculum of BS
%  Information Technology, p. 65), whose course outline is word for
%  word identical to HEC 2017 (p. 89).  Computing Core, compulsory,
%  semester 7, 3 credit hours, prerequisite Data Structures.
%
%  The preamble is shared with the sibling courses' schedule pages.
%
%  Build:  pdflatex AoA_Fall2026_Lecture_Schedule.tex
%=====================================================================
\documentclass[10pt,a4paper]{article}
\usepackage[a4paper,left=1.2cm,right=1.2cm,top=0.85cm,bottom=0.6cm]{geometry}
\usepackage[T1]{fontenc}
\usepackage{lmodern}
\usepackage{microtype}
\usepackage[table]{xcolor}
\usepackage{array,ragged2e}
\usepackage{graphicx}
\usepackage{fontawesome5}
\pagestyle{empty}
\pdfinfo{/Title (Analysis of Algorithms, Fall 2026: Weekly Lecture Schedule)}
\setlength{\parindent}{0pt}
\setlength{\parskip}{0pt}

%--- palette: shared BSIT palette (../Data Science Fall 2026/dshandout.sty
%    family); part colours darkened a little so white chip text stays
%    legible at 6 pt
\definecolor{primary}{HTML}{1A5276}
\definecolor{secondary}{HTML}{2E86C1}
\definecolor{accent}{HTML}{C0392B}
\definecolor{lightbg}{HTML}{EBF5FB}
\definecolor{neutral}{HTML}{707B7C}
\definecolor{sub}{HTML}{465152}
\definecolor{p1}{HTML}{1A5276}
\definecolor{p2}{HTML}{2874A6}
\definecolor{p3}{HTML}{1E8449}
\definecolor{p4}{HTML}{8E44AD}
\definecolor{p5}{HTML}{117A65}
\definecolor{p6}{HTML}{8A6D0A}

%--- building blocks -------------------------------------------------
\newcommand{\chip}[2]{{\setlength{\fboxsep}{1.3pt}\colorbox{#1}{\makebox[1.45em][c]{%
  \fontsize{6.4}{7}\selectfont\bfseries\color{white}#2}}}}
\newcommand{\chipw}[2]{{\setlength{\fboxsep}{1.3pt}\colorbox{#1}{%
  \fontsize{6.4}{7}\selectfont\bfseries\color{white}\,#2\,}}}
% \ch{part}{number}{title}{what it covers}  -- one chapter, own paragraph
\newcommand{\ch}[4]{\chip{p#1}{#2}\,\textbf{#3}\ {\color{neutral}--}\ {\fontsize{7.3}{8.5}\selectfont\color{sub}#4}\par}
\newcommand{\ex}[1]{{\bfseries\color{accent}#1}}      % exams and the big deliverables
\newcommand{\ev}[1]{{\bfseries\color{primary}#1}}     % quizzes, assignments, project points
\newcommand{\none}{{\color{neutral}---}}
\newcommand{\hdr}[1]{{\bfseries\color{white}#1}}
\newcommand{\lab}[1]{{\color{sub}#1}}

\newlength{\topicw}   % set inside the document, once \linewidth and \tabcolsep are final
\newcolumntype{K}{>{\centering\arraybackslash}p{0.75cm}}
\newcolumntype{T}{>{\RaggedRight\arraybackslash}p{\topicw}}
\newcolumntype{L}[1]{>{\RaggedRight\arraybackslash}p{#1}}
% \leavevmode: a cell that opens with \color puts a whatsit before its paragraph, and
% \vtop then aligns it by its top edge -- one strut lower than a cell that opens with text
\newcommand{\wk}[4]{\leavevmode{\large\bfseries\color{primary}#1} & \leavevmode #2 &
  \leavevmode #3 & \leavevmode #4 \\ \hline}

\newcommand{\fact}[3]{{\setlength{\fboxsep}{4pt}\colorbox{lightbg}{%
  \parbox[t]{\dimexpr#1-8pt\relax}{\RaggedRight\fontsize{6.2}{7}\selectfont\bfseries\color{neutral}\MakeUppercase{#2}\\[1.5pt]%
  \fontsize{7.8}{9.2}\selectfont\bfseries\color{primary}#3}}}}
\newcommand{\tile}[3]{{\colorbox{lightbg}{\parbox[c][0.9cm][c]{\dimexpr(\linewidth-3pt*\numexpr#1-1\relax)/#1-6pt-0.5pt\relax}{%
  \centering{\large\bfseries\color{primary}#2}\\[-0.5pt]{\fontsize{6.6}{7.4}\selectfont\color{sub}#3}}}}}

% \legend{I}{II}{III}{IV}{V}{VI} -- the six part names, as a 3 x 2 grid of colour chips
\newcommand{\legend}[6]{{\fontsize{6.6}{8}\selectfont\color{sub}%
  \begin{tabular}{@{}p{0.31\linewidth}p{0.31\linewidth}p{0.31\linewidth}@{}}
  \chipw{p1}{I}\ #1 & \chipw{p2}{II}\ #2 & \chipw{p3}{III}\ #3 \\[1pt]
  \chipw{p4}{IV}\ #4 & \chipw{p5}{V}\ #5 & \chipw{p6}{VI}\ #6
  \end{tabular}}}

\begin{document}
\fontsize{8.2}{9.8}\selectfont

%--- header ----------------------------------------------------------
\begin{minipage}[c]{1.5cm}%
  \IfFileExists{assets/iub-logo.png}{\includegraphics[height=1.15cm]{assets/iub-logo.png}}{}%
\end{minipage}\hfill
\begin{minipage}[c]{0.65\linewidth}\centering
  {\footnotesize\bfseries\color{primary}THE ISLAMIA UNIVERSITY OF BAHAWALPUR}\\[1pt]
  {\scriptsize\color{neutral}Department of Information Technology}
\end{minipage}\hfill
\begin{minipage}[c]{1.5cm}\raggedleft
  \IfFileExists{assets/dit-logo.png}{\includegraphics[height=1.15cm]{assets/dit-logo.png}}{}%
\end{minipage}

\vspace{3pt}
{\color{secondary}\rule{\linewidth}{1.4pt}}\\[2pt]
{\Large\bfseries\color{primary}ANALYSIS OF ALGORITHMS}\hfill
{\normalsize\bfseries\color{secondary}Weekly Lecture Schedule\ \textbullet\ Fall 2026}\\[1pt]
{\footnotesize\color{sub}Bachelor of Science in Information Technology\ \textbullet\
 \emph{Design, Proof and Measurement}}\\[1pt]
{\footnotesize\color{sub}HEC 2023 Computing Core, semester 7. Every bound in this course is proved, and then measured on a real machine.}\\[1pt]
{\color{secondary}\rule{\linewidth}{0.6pt}}

\vspace{4pt}
\fact{0.235\linewidth}{Credit hours}{3 (3--0): 3 lecture hours a week}\hfill
\fact{0.235\linewidth}{Duration}{16 teaching weeks; midterm Week 8, project due Week 16}\hfill
\fact{0.235\linewidth}{Prerequisite}{Data Structures (HEC 2023, p.\ 65)}\hfill
\fact{0.235\linewidth}{Contents}{23 chapters in 6 parts; 23 measured experiments in C++}

\vspace{5pt}
%--- the schedule ----------------------------------------------------
\begingroup
\renewcommand{\arraystretch}{1.3}
\setlength{\tabcolsep}{3.5pt}
\setlength{\topicw}{\dimexpr\linewidth-0.75cm-4.05cm-3.3cm-6\tabcolsep\relax}
\arrayrulecolor{primary!18}
\rowcolors{2}{white}{lightbg}
\noindent\begin{tabular}{@{}K T L{4.05cm} L{3.3cm}@{}}
\rowcolor{primary}
\hdr{\faIcon{calendar-alt}} & \hdr{\faIcon{book-open}\ \ Lecture topics} &
\hdr{\faIcon{flask}\ \ Measured experiment} & \hdr{\faIcon{clipboard-check}\ \ Assessment} \\
\wk{1}
 {\ch{1}{1}{The Role of Algorithms}{what an algorithm is; the unit-cost RAM model and the three things it ignores; why the question is asymptotic}
  \ch{1}{2}{Input Size, Input Nature, Three Cases}{size as bit length; best, worst and average case; why an average needs a stated distribution}}
 {\lab{The ratio test; the crossover at $n \approx 65$}}{\none}
\wk{2}
 {\ch{1}{3}{Asymptotic Notation}{$O$, $\Omega$, $\Theta$, $o$ and $\omega$ defined and proved from the definitions; the limit test; four common abuses}}
 {\lab{Fit a curve; exclude the exponents either side}}{\none}
\wk{3}
 {\ch{1}{4}{Correctness: Loop Invariants}{initialisation, maintenance, termination; partial against total correctness; the loop variant}}
 {\lab{A bug that is never wrong and never finishes}}{\none}
\wk{4}
 {\ch{2}{5}{Recurrences I: Substitution and Trees}{guess and prove by induction; strengthening the hypothesis; recursion trees; unbalanced splits}
  \ch{2}{6}{Recurrences II: The Master Theorem}{all three cases, the regularity condition, and a recurrence that falls in the gap}}
 {\lab{Recursion trees counted; the theorem verified exactly}}{\none}
\wk{5}
 {\ch{2}{7}{Brute Force, and When It Is Right}{exhaustive search as baseline and oracle; maximum subarray in $n^3$, $n^2$ and $n$}
  \ch{2}{8}{Divide and Conquer: Merge Sort}{the three-step pattern; correctness from the merge invariant; tuning the cutoff}}
 {\lab{Where $n^2$ beats $n \lg n$; merge sort tuned}}{\none}
\wk{6}
 {\ch{2}{9}{Quicksort and Randomisation}{the worst case is already-sorted input; expected $\Theta(n \lg n)$ on every input; introsort}}
 {\lab{The adversarial input, and what randomising costs}}{\ev{Quiz 1} (Ch 1--6)}
\wk{7}
 {\ch{3}{10}{Lower Bounds and Linear-Time Sorting}{the decision-tree $\Omega(n \lg n)$ proof; counting and radix sort, and what each one assumes}}
 {\lab{Counting and radix sort against \texttt{std::sort}}}{\ev{Assignment 1 set}}
\wk{8}
 {\ch{3}{11}{Selection and Order Statistics}{quickselect; why deleting one recursive call changes the master-theorem case; median of medians}}
 {\lab{Midterm, then quickselect against a full sort}}{\ex{Midterm} (Ch 1--9)}
\wk{9}
 {\ch{3}{12}{Heaps and Search Trees, Analysed}{why a heap's height cannot fail and a search tree's can; Build-Heap is $\Theta(n)$}
  \ch{3}{13}{Hashing, and What $O(1)$ Promises}{simple uniform hashing; load factor; the two ways the assumption fails}}
 {\lab{A hash table degrading under a poor hash}}{\ev{Assignment 1 due}}
\wk{10}
 {\ch{4}{14}{Greedy Algorithms}{the greedy-choice property; exchange arguments; activity selection; Huffman coding proved optimal}}
 {\lab{Huffman coding against the entropy bound}}{\none}
\wk{11}
 {\ch{4}{15}{Dynamic Programming I: The Elements}{optimal substructure and overlapping subproblems; memoisation against bottom-up}
  \ch{4}{16}{Dynamic Programming II: The Classics}{LCS, edit distance, 0/1 knapsack; why $\Theta(nW)$ is pseudo-polynomial}}
 {\lab{Memoised against bottom-up, timed and counted}}{\none}
\wk{12}
 {\ch{4}{17}{Amortised Analysis}{aggregate, accounting and potential methods; why $n$ appends to a doubling array cost $\Theta(n)$}}
 {\lab{Aggregate cost against worst-case cost}}{\ev{Quiz 2} (Ch 10--16)}
\wk{13}
 {\ch{5}{18}{Graphs: Representation and Sparsity}{matrix against list; why every bound carries two parameters; BFS, DFS, topological order}
  \ch{5}{19}{Minimum Spanning Trees}{the cut property proved; Kruskal and Prim; union--find, where amortised analysis pays off}}
 {\lab{The matrix/list crossover as density falls}}{\ev{Assignment 2 set}\newline \ev{Project proposal}}
\wk{14}
 {\ch{5}{20}{Shortest Paths}{Dijkstra with its correctness proof; Bellman--Ford and negative edges; Floyd--Warshall for all pairs}}
 {\lab{Dijkstra with and without a priority queue}}{\none}
\wk{15}
 {\ch{6}{21}{String Matching}{naive, Rabin--Karp and KMP; the prefix function; why KMP has no bad case at all}
  \ch{6}{22}{Complexity Classes: P, NP, NP-Complete}{verifying against finding; reductions; what to do when a problem is NP-complete}}
 {\lab{Naive, Rabin--Karp and KMP on one corpus}}{\ev{Quiz 3} (Ch 17--22)\newline \ev{Assignment 2 due}}
\wk{16}
 {\ch{6}{23}{The Semester Project}{implement, measure and defend; the ten rules of experimental method; the measurement harness}}
 {\lab{Implement, measure, defend}}{\ex{Project due; Final} (Ch 1--22)}
\end{tabular}
\endgroup

\vspace{4pt}
%--- assessment weights ----------------------------------------------
{\fontsize{6.6}{8}\selectfont\bfseries\color{neutral}\MakeUppercase{Assessment weights}}\\[2pt]
\tile{3}{20\%}{Sessional: quizzes, assignments, project}\hfill
\tile{3}{30\%}{Mid-term, Week 8}\hfill
\tile{3}{50\%}{Final-term, all chapters}

\vspace{3pt}
%--- legend and notes ------------------------------------------------
\legend{Analysing an Algorithm}{Recursion and Divide-and-Conquer}{Sorting, Selection and the Structures Beneath}
       {The Design Paradigms}{Graph Algorithms}{Strings, Limits and Practice}

\vspace{1pt}
{\fontsize{6.3}{7.6}\selectfont\color{neutral}%
Follows HEC 2023 (\emph{Curriculum of BS Information Technology}, p.\ 65), whose course outline is word for word identical to HEC 2017 (p.\ 89). Three topics are taught that the outline does not name, each justified in the handout's Course Specification: the comparison lower bound (Week 7), the first result in the course about a \emph{problem} rather than an algorithm; minimum spanning trees (Week 13), absent from the outline text but required by CLO-7 by name; and amortised analysis (Week 12), without which Chapter 13's hash-table claim and Chapter 19's union--find cannot be stated honestly, and which CLO-3 requires for sequences of operations. The published reading list misspells Cormen as ``Corman'' and cites CLRS 3rd ed.\ (2009); Appendix~E of the handout cites the 4th (2022). Chapter numbers refer to the handout, \emph{Analysis of Algorithms --- Design, Proof and Measurement} (242 pp., 23 chapters, 7 appendices); Appendix~G of that handout is the source of record for this page. Every chapter carries one experiment whose figures were produced on the machine described in Appendix~A and pasted into the text; the sources are in \texttt{code/} and each is reproducible.
Weights follow the department's standard split: Sessional 20\%, Mid-term 30\%, Final-term 50\%.\par}

\end{document}
